\documentclass[10pt]{article}
\usepackage[T1]{fontenc}
\usepackage[utf8]{inputenc}
\usepackage[french]{babel}
\usepackage{amsmath,amssymb,mathrsfs,mathpazo}
\usepackage{geometry}
\geometry{a4paper,margin=2cm}
\usepackage{graphicx,array,tabularx,tabularray,colortbl,multirow,multicol,booktabs,enumitem,pifont,ulem,fancybox,diagbox,cancel}
\usepackage{pstricks,pstricks-add,pst-plot,pst-tree,pst-node,pst-text,pst-func,pst-math,pst-bspline,pst-3dplot,pst-eucl,pst-solides3d}
\let\pscalculate\undefined
\usepackage{pst-bezier}
\usepackage{tikz,pgfplots,tkz-tab}
\usetikzlibrary{arrows,arrows.meta,patterns,calc,positioning,trees,decorations.markings,shapes.geometric,shapes.misc,shapes.arrows,intersections,decorations.pathmorphing,decorations.pathreplacing,backgrounds,shadows,babel}
\pgfplotsset{compat=1.18}
\usepgfplotslibrary{fillbetween}
\usepackage[np]{numprint}
\usepackage{listings,makecell,pas-tableur,twoopt}
\let\barre\undefined
\usepackage[pstricks]{tablvar}
\definecolor{darkgreen}{RGB}{0,100,0}
\providecommand{\red}{\color{red}}
\providecommand{\blue}{\color{blue}}
\providecommand{\green}{\color{green}}
\providecommand{\black}{\color{black}}
\providecommand{\white}{\color{white}}
\providecommand{\gray}{\color{gray}}
\providecommand{\cyan}{\color{cyan}}
\providecommand{\magenta}{\color{magenta}}
\providecommand{\yellow}{\color{yellow}}
\providecommand{\violet}{\color{violet}}
\providecommand{\orange}{\color{orange}}
\providecommand{\euro}{\texteuro}
\providecommand{\EUR}{\texteuro}
\providecommand{\ieme}{\textsuperscript{e}}
\providecommand{\degres}{\ensuremath{^\circ}}
\providecommand{\degree}{\ensuremath{^\circ}}
\providecommand{\vv}[1]{\overrightarrow{#1}}
\providecommand{\Coord}[1]{#1}
\providecommand{\decosix}{}
\providecommand{\starredbullet}{\ensuremath{\bullet}}
\pagestyle{empty}
\setlength{\parindent}{0pt}
\usepackage{adjustbox}\DeclareUnicodeCharacter{2236}{:}\DeclareUnicodeCharacter{2190}{\ensuremath{\leftarrow}}
\providecommand{\R}{\mathbb{R}}
\providecommand{\N}{\mathbb{N}}
\providecommand{\Z}{\mathbb{Z}}
\providecommand{\Q}{\mathbb{Q}}
\providecommand{\D}{\mathbb{D}}
\providecommand{\C}{\mathbb{C}}
\providecommand{\e}{{\mathrm{e}}}
\providecommand{\vect}[1]{\overrightarrow{#1}}
\providecommand{\vectt}[1]{\overrightarrow{\mathrm{#1}}}
\providecommand{\vvt}[1]{\overrightarrow{\mathrm{#1}}}
\providecommand{\np}[2][]{{#2~#1}}
\providecommand{\vv}[1]{{\overrightarrow{#1}}}
\renewcommand{\d}{{\mathrm{d}}}
\providecommand{\degres}{{^\circ}}
\providecommand{\degre}{{^\circ}}
\providecommand{\blue}{\color{blue}}
\providecommand{\red}{\color{red}}
\providecommand{\enskip}{\quad}
\providecommand{\llbracket}{\lbrack\!\lbrack}
\providecommand{\rrbracket}{\rbrack\!\rbrack}
\providecommand{\euro}{{\text{€}}}
\providecommand{\no}{{\text{n°}}}
\providecommand{\interval}[2]{{[#1;#2]}}
\providecommand{\intervFO}[2]{{[#1;#2[}}
\providecommand{\dotfill}{\ldots}

\providecommand{\ds}{\displaystyle}
\providecommand{\textding}[1]{#1}
\providecommand{\Oij}{\(\mathrm{O};\vec{\imath},\vec{\jmath}\)}
\providecommand{\Oijk}{\(\mathrm{O};\vec{\imath},\vec{\jmath},\vec{k}\)}
\providecommand{\Ouv}{\(\mathrm{O};\vec{u},\vec{v}\)}
\providecommand{\intervFF}[2]{\ensuremath{[#1;#2]}}
\renewcommand{\euro}{\ifmmode\text{\texteuro}\else\texteuro\fi}
\DeclareUnicodeCharacter{00B0}{\ifmmode^{\circ}\else\textdegree\fi}
\DeclareUnicodeCharacter{20AC}{\ifmmode\text{\texteuro}\else\texteuro\fi}
% Énoncés : Éducation nationale. Transcription et corrigés : APMEP, Denis Vergès et contributeurs. Découpage et classement : CoopMaths, Sébastien Lozano, Éric Elter et contributeurs. Adaptation de la présentation : Maths Explorer.
% CC BY-SA : https://coopmaths.fr/www/about/
% https://coopmaths.fr/alea/static/bac/2026/tex/bac_2026_05_sujet1_ameriquenord_4.tex
\begin{document}
\noindent\textit{APMEP / CoopMaths. CC BY-SA. Présentation adaptée par Maths Explorer.}\par\medskip
On considère la fonction $f$ définie sur $\R$ par 

\[f(x)=5 \ln \left(x^{2}+1\right)-3 x\]
et on admet que la fonction $f$ est dérivable sur $\R$.

On note $\mathcal{C}_{f}$ sa courbe représentative dans un repère orthonormé du plan.

On a tracé ci-dessous la courbe $\mathcal{C}_{f}$ et la tangente à la courbe $\mathcal{C}_{f}$ au point A d'abscisse 1.

\begin{center}
\begin{tikzpicture}
\begin{axis}[
	x={9mm}, y={9mm},
	xmin = -6.5, xmax = 6.9,
	ymin = -2.7, ymax = 6.1,
	xtick distance=1,
	ytick distance=1,
	minor tick num = 0,
	grid=none, axis lines =center]
	\node[below left=3pt] at (axis cs: 0,0) {{\footnotesize 0}};
	\addplot [line width=1.2pt, color=red, smooth, samples=100, domain= -1.5:6.9 ]{5*ln(x*x+1)-3*x};
	\node[red] at(axis cs: -1,3) {$\mathcal{C}_f$};
	\addplot [line width=1pt, color=teal, smooth, samples=10, domain= -1.5:6.9 ]{2*x-1.534};
	\draw[shift={(axis cs: 1,0.466)}] (-2pt,-2pt)--(2pt,2pt) (-2pt,2pt)node[above left]{A}--(2pt,-2pt) ;
\end{axis}
\end{tikzpicture}
\end{center}













\begin{enumerate}
\item Conjecturer, à l'aide de la représentation graphique de la fonction $f$, les intervalles de $\R$ sur lesquels la fonction $f$ semble convexe ou concave.
\item Déterminer, en justifiant, la limite de la fonction $f$ en $-\infty$.
\item
	\begin{enumerate}
		\item Démontrer que, pour tout $x$ réel strictement positif,

\[f(x)=x\left(10 \dfrac{\ln x}{x}-3\right)+5 \ln \left(1+\dfrac{1}{x^{2}}\right)\]

		\item Déterminer, en justifiant, la limite de la fonction $f$ en $+\infty$.\end{enumerate}
\item
	\begin{enumerate}
		\item Démontrer que pour tout $x$ réel, \quad $f'(x)= \dfrac{-3 x^{2}+10 x-3}{x^{2}+1}$.
		\item Étudier les variations de la fonction $f$ sur $\R$.
	\end{enumerate}
	\item On admet que la fonction $f$ est deux fois dérivable sur $\R$ et que pour tout réel $x$,
	\[f''(x)=\dfrac{-10 x^{2}+10}{\left(x^{2}+1\right)^{2}}\]

	\begin{enumerate}
		\item Valider ou rejeter la conjecture faite à la question \textbf{1}.
		\item Déterminer l'équation réduite de la tangente à la courbe $\mathcal{C}_{f}$ au point A d'abscisse 1.
		\item En déduire que pour tout $x \geqslant 1$,\quad $\ln \left(x^{2}+1\right) \leqslant x+\ln (2)-1$.
	\end{enumerate}
	\end{enumerate}

\end{document}
